\documentclass[11pt]{article}
\usepackage[margin=1in]{geometry}
\usepackage[T1]{fontenc}
\usepackage[utf8]{inputenc}
\usepackage{lmodern}
\usepackage{setspace}
\usepackage{booktabs}
\usepackage{longtable}
\usepackage{array}
\usepackage{amsmath}
\usepackage{amssymb}
\usepackage{hyperref}
\usepackage{xcolor}

\hypersetup{
    colorlinks=true,
    linkcolor=blue,
    urlcolor=blue,
    citecolor=blue
}

\setlength{\parskip}{0.5em}
\setlength{\parindent}{0pt}
\onehalfspacing

\title{Hexapod Brain Roadmap\\\large Whiteboard Strategy Integrated with Current Firmware}
\author{E90 Hexapod Team}
\date{\today}

\begin{document}
\maketitle

\section{Purpose}

This document turns the professor meeting whiteboard plan into a concrete implementation roadmap grounded in the current driving code:

\begin{itemize}
    \item \texttt{adaptive\_ondevice\_hexapod/firmware/hexapod\_walking\_v2/hexapod\_walking\_v2.ino}
\end{itemize}

The goal is to evolve from a deterministic tripod gait controller into an adaptive ``brain'' that remains safe on hardware while progressively introducing reflexes, online adaptation, and learning-based policy support.

\section{Current System Baseline (What Already Exists)}

The current firmware already provides a strong and safe base:

\begin{enumerate}
    \item \textbf{12-DOF servo actuation with calibration}
    \begin{itemize}
        \item Explicit channel map for all 6 legs and 2 joints per leg (abduction/flexion)
        \item Per-servo direction and center calibration
    \end{itemize}
    \item \textbf{Reliable actuator interface}
    \begin{itemize}
        \item Pololu Maestro serial protocol
        \item Hard pulse clamps (\texttt{MIN\_US}, \texttt{MAX\_US})
    \end{itemize}
    \item \textbf{Deterministic tripod gait state machine}
    \begin{itemize}
        \item Tripod A: lift $\rightarrow$ swing+push $\rightarrow$ plant
        \item Tripod B: lift $\rightarrow$ swing+push $\rightarrow$ plant
    \end{itemize}
    \item \textbf{Smooth interpolation and phase timing}
    \begin{itemize}
        \item \texttt{moveSmoothTo(...)} with independent step counts and delays
        \item Existing tunable gait constants (stride amplitudes, flex targets, phase durations)
    \end{itemize}
\end{enumerate}

\textbf{Interpretation:} the robot already has a practical ``brain v0.'' The roadmap should preserve this safety backbone and layer adaptation on top rather than replacing it.

\section{Whiteboard Ideas Mapped to Firmware Reality}

\begin{table}[h]
\centering
\begin{tabular}{>{\raggedright\arraybackslash}p{0.30\linewidth} >{\raggedright\arraybackslash}p{0.62\linewidth}}
\toprule
\textbf{Whiteboard concept} & \textbf{Concrete mapping in current stack} \\
\midrule
TA/TB clock phases & Existing Tripod A/Tripod B phase sequence in firmware loop \\
P0/P1/P2/P3 phase points & Keyframe targets built through \texttt{buildPose(...)} per sub-phase \\
Waveform-based generation & Optional future gait generator mode; keep keyframe mode as default safety path \\
PPO policy block & Learn bounded updates to gait parameters (not raw servo microseconds) \\
EOM/predictive block & Short-horizon stability predictor from IMU + gait phase for shaping/safety \\
Terrain/speed conditioning & Add condition vector (terrain class + speed target) to select parameter presets or policy outputs \\
\bottomrule
\end{tabular}
\caption{Converting high-level meeting concepts into implementation-ready modules}
\end{table}

\section{Target Architecture}

To stay robust on hardware, architecture is layered and conservative:

\subsection*{Layer 1: Execution (existing)}
\begin{itemize}
    \item Servo command generation, clamping, and Maestro output
    \item Interpolated motion between pose targets
\end{itemize}

\subsection*{Layer 2: Parameterized gait generator (next)}
\begin{itemize}
    \item Convert fixed constants into runtime state
    \item Preserve tripod keyframes as canonical gait
\end{itemize}

\subsection*{Layer 3: Fast reflex loop (next)}
\begin{itemize}
    \item IMU-based roll/pitch and angular-rate stabilization
    \item Applied every interpolation tick
\end{itemize}

\subsection*{Layer 4: Slow adaptation loop (next)}
\begin{itemize}
    \item Once per gait cycle/window, update baseline parameters
    \item Enforce strict safety bounds and slew-rate limits
\end{itemize}

\subsection*{Layer 5: Learning policy (later)}
\begin{itemize}
    \item PPO trained in simulation with domain randomization
    \item Output high-level bounded parameters only
\end{itemize}

\subsection*{Layer 6: Predictor (optional/late)}
\begin{itemize}
    \item Estimate short-horizon instability risk
    \item Use for reward shaping and cautious-mode trigger
\end{itemize}

\section{Leg Coupling Specification (Concrete)}

\subsection*{Leg indexing and tripod groups}

Define leg set:
\[
\mathcal{L}=\{\mathrm{FL},\mathrm{FR},\mathrm{ML},\mathrm{MR},\mathrm{RL},\mathrm{RR}\}
\]

Tripod grouping (matching current firmware):
\[
\mathcal{A}=\{\mathrm{FL},\mathrm{MR},\mathrm{RL}\},\quad
\mathcal{B}=\{\mathrm{FR},\mathrm{ML},\mathrm{RR}\}
\]

Joint channels per leg:
\[
q_l = (q^{\mathrm{abd}}_l,\ q^{\mathrm{flx}}_l)
\]

\subsection*{Phase-coupled clock}

Use normalized gait clock $\varphi\in[0,1)$ over one full stride.
Tripod phase offsets:
\[
\delta_l=
\begin{cases}
0.0, & l\in\mathcal{A}\\
0.5, & l\in\mathcal{B}
\end{cases}
\]
and per-leg local phase:
\[
\varphi_l = (\varphi + \delta_l)\bmod 1
\]

This enforces anti-phase coupling between tripods.

\subsection*{Sub-phase schedule per leg}

Each leg has three sub-phases using local phase $\varphi_l$:
\begin{itemize}
    \item Lift: $\varphi_l\in[0,\tau_1)$
    \item Swing: $\varphi_l\in[\tau_1,\tau_2)$
    \item Plant/Push: $\varphi_l\in[\tau_2,1)$
\end{itemize}

Initial values from v2 timing ratios:
\[
\tau_1 \approx \frac{\texttt{LIFT\_STEPS}}{\texttt{LIFT\_STEPS}+\texttt{SWING\_STEPS}+\texttt{PLANT\_STEPS}+\texttt{PUSH\_STEPS}},
\]
\[
\tau_2 \approx \frac{\texttt{LIFT\_STEPS}+\texttt{SWING\_STEPS}}{\texttt{LIFT\_STEPS}+\texttt{SWING\_STEPS}+\texttt{PLANT\_STEPS}+\texttt{PUSH\_STEPS}}
\]

\subsection*{Per-leg command law}

Let $b_l$ be baseline abduction offset (front/mid/rear specific), and $a_l$ be stride amplitude (front/mid/rear specific). Define signed stride direction:
\[
\sigma_l(\varphi_l)=
\begin{cases}
+1, & \text{swing}\\
-1, & \text{push}\\
0,  & \text{lift/plant transition}
\end{cases}
\]

Abduction command:
\[
q^{\mathrm{abd}}_l = b_l + s_{\mathrm{stride}}\cdot a_l\cdot \sigma_l(\varphi_l)
\]

Flexion command:
\[
q^{\mathrm{flx}}_l=
\begin{cases}
q_{\mathrm{lift}},  & \text{lift}\\
q_{\mathrm{lift}},  & \text{swing}\\
q_{\mathrm{plant}}, & \text{plant transition}\\
q_{\mathrm{push}},  & \text{stance push}
\end{cases}
\]

where:
\[
q_{\mathrm{lift}} = q^0_{\mathrm{lift}}\cdot s_{\mathrm{height}},\quad
q_{\mathrm{push}} = q^0_{\mathrm{push}},\quad
q_{\mathrm{plant}} = q^0_{\mathrm{start}}
\]

\subsection*{Roll/pitch coupling distribution across legs}

Reflex corrections are coupled by side and body-axis:
\begin{align}
\Delta_{\mathrm{roll}} &= K_{\phi}\phi + K_{\dot{\phi}}\dot{\phi}\\
\Delta_{\mathrm{pitch}} &= K_{\theta}\theta + K_{\dot{\theta}}\dot{\theta}
\end{align}

For each leg $l$, define side sign $s_l\in\{+1,-1\}$ and front-rear sign $r_l\in\{+1,0,-1\}$:
\begin{itemize}
    \item $s_l=+1$ for left legs (FL, ML, RL), $s_l=-1$ for right legs (FR, MR, RR)
    \item $r_l=-1$ front (FL, FR), $r_l=0$ middle (ML, MR), $r_l=+1$ rear (RL, RR)
\end{itemize}

Then leg flex bias is:
\[
\Delta q^{\mathrm{flx}}_l = s_l\Delta_{\mathrm{roll}} + r_l\Delta_{\mathrm{pitch}}
\]

Final flex command before pulse conversion:
\[
\hat{q}^{\mathrm{flx}}_l = q^{\mathrm{flx}}_l + \Delta q^{\mathrm{flx}}_l
\]
followed by hard clamping and slew-rate limiting.

\subsection*{Why this coupling is implementable now}

This coupling model is directly compatible with the existing v2 structure:
\begin{itemize}
    \item Current tripod functions already encode anti-phase A/B behavior
    \item Existing front/mid/rear amplitudes map to $a_l$
    \item Existing start/lift/plant/push flex constants map to $q^0_{\cdot}$
    \item Existing clamp and interpolation functions enforce safety and smoothness
\end{itemize}

\section{Core Parameterization}

\subsection*{Runtime gait parameter vector}

Define the control parameter vector:
\[
\mathbf{p}=
\begin{bmatrix}
q^0_{\mathrm{start}} & q^0_{\mathrm{push}} & q^0_{\mathrm{lift}} & q^0_{\mathrm{plant}} &
a_F & a_M & a_R & \omega &
s_{\mathrm{stride}} & s_{\mathrm{height}}
\end{bmatrix}^{\top}
\]
where $(a_F,a_M,a_R)$ are front/middle/rear abduction amplitudes and $\omega$ is gait angular frequency.

\subsection*{State, action, and observation}

At fast tick $k$ (period $\Delta t$), define:
\[
\mathbf{x}_k=
\begin{bmatrix}
\phi_k & \theta_k & \dot{\phi}_k & \dot{\theta}_k & \varphi_k
\end{bmatrix}^{\top},\quad
\mathbf{u}_k=
\begin{bmatrix}
\Delta \mathbf{p}_k
\end{bmatrix}
\]
with observation $\mathbf{o}_k=\mathbf{x}_k$ in the minimal onboard setup.

\subsection*{Online adaptation parameter subset}

For stability-first adaptation on hardware:
\[
\boldsymbol{\theta}=
\begin{bmatrix}
K_{\phi} & K_{\dot{\phi}} & K_{\theta} & K_{\dot{\theta}} &
s_{\mathrm{stride}} & s_{\mathrm{height}} & s_{\mathrm{freq}}
\end{bmatrix}^{\top}
\]
This is the vector updated online (SPSA/bandit/PPO output), while core safety bounds remain fixed.

\section{Control and Stability Formulation}

\subsection*{Phase oscillator and hybrid gait dynamics}

Global gait phase:
\begin{equation}
\varphi_{k+1} = (\varphi_k + \omega_k \Delta t)\bmod 1
\end{equation}

Per-leg local phase:
\begin{equation}
\varphi_{l,k} = (\varphi_k + \delta_l)\bmod 1,\quad
\delta_l\in\{0,0.5\}
\end{equation}

Piecewise gait mode map:
\[
m(\varphi_{l,k})\in\{\text{lift},\text{swing},\text{plant},\text{push}\}
\]
with switching boundaries $(\tau_1,\tau_2,\tau_3)$ obtained from step-count ratios.

\subsection*{Per-leg command generator}

For each leg $l$:
\begin{align}
q^{\mathrm{abd}}_{l,k} &= b_l + s_{\mathrm{stride},k}\,a_{g(l)}\,\sigma(m(\varphi_{l,k}))\\
q^{\mathrm{flx,base}}_{l,k} &=
\begin{cases}
q^0_{\mathrm{lift}}\,s_{\mathrm{height},k}, & m=\text{lift/swing}\\
q^0_{\mathrm{plant}}, & m=\text{plant}\\
q^0_{\mathrm{push}}, & m=\text{push}
\end{cases}
\end{align}
where $g(l)\in\{F,M,R\}$ maps each leg to front/middle/rear class.

\subsection*{Reflex-coupled stabilization law}

\begin{align}
\Delta_{\mathrm{roll},k} &= K_{\phi}\phi_k + K_{\dot{\phi}}\dot{\phi}_k\\
\Delta_{\mathrm{pitch},k} &= K_{\theta}\theta_k + K_{\dot{\theta}}\dot{\theta}_k
\end{align}

Leg-specific flex coupling:
\begin{equation}
\Delta q^{\mathrm{flx}}_{l,k}=s_l\Delta_{\mathrm{roll},k}+r_l\Delta_{\mathrm{pitch},k}
\end{equation}
with side sign $s_l\in\{+1,-1\}$ and front-rear sign $r_l\in\{-1,0,+1\}$.

Final command before pulse conversion:
\begin{equation}
\hat{q}^{\mathrm{flx}}_{l,k}=q^{\mathrm{flx,base}}_{l,k}+\Delta q^{\mathrm{flx}}_{l,k}
\end{equation}

\subsection*{Safety projection and slew constraints}

Pulse command for servo channel $c(l,j)$:
\begin{equation}
u^{\mu s}_{c,k}=\mathrm{clip}\!\left(\mu_c+d_c\hat{q}_{l,j,k},\ u^{\min}_c,\ u^{\max}_c\right)
\end{equation}
where $\mu_c$ is servo center and $d_c\in\{-1,+1\}$ is direction.

Slew-limited command:
\begin{equation}
u^{\mu s}_{c,k}\leftarrow
\mathrm{clip}\!\left(u^{\mu s}_{c,k},\ u^{\mu s}_{c,k-1}-\Delta^{\max}_c,\ u^{\mu s}_{c,k-1}+\Delta^{\max}_c\right)
\end{equation}

\subsection*{Instability metric and mode switching}

Instantaneous score:
\begin{equation}
S_k=w_1|\phi_k|+w_2|\theta_k|+w_3|\dot{\phi}_k|+w_4|\dot{\theta}_k|
\end{equation}
Windowed RMS:
\begin{equation}
S_{\mathrm{rms}}(n)=\sqrt{\frac{1}{N}\sum_{i=n-N+1}^{n}S_i^2}
\end{equation}

Hybrid mode:
\[
\texttt{mode}=
\begin{cases}
\texttt{cautious}, & S_{\mathrm{rms}}>\gamma_H \text{ for } M_H \text{ cycles}\\
\texttt{nominal}, & S_{\mathrm{rms}}<\gamma_L \text{ for } M_L \text{ cycles}
\end{cases}
\]
with $\gamma_H>\gamma_L$ (hysteresis).

\subsection*{Online adaptation objective and update}

Per-window cost:
\begin{equation}
J=
\alpha_1\,\mathrm{RMS}(\phi)+
\alpha_2\,\mathrm{RMS}(\theta)+
\alpha_3\,\mathrm{mean}(|\dot{\phi}|+|\dot{\theta}|)+
\alpha_4\,\mathbb{I}_{\mathrm{fall}}+
\alpha_5\,E_{\mathrm{cmd}}
\end{equation}
where
\[
E_{\mathrm{cmd}}=\frac{1}{NT}\sum_{k=1}^{N}\sum_{c=1}^{12}|u^{\mu s}_{c,k}-u^{\mu s}_{c,k-1}|
\]
is a command-variation proxy for effort.

SPSA update:
\begin{align}
\boldsymbol{\theta}^{+}&=\boldsymbol{\theta}_n+c_n\boldsymbol{\Delta}_n,\quad
\boldsymbol{\theta}^{-}=\boldsymbol{\theta}_n-c_n\boldsymbol{\Delta}_n\\
\hat{\mathbf{g}}_n&=\frac{J(\boldsymbol{\theta}^{+})-J(\boldsymbol{\theta}^{-})}{2c_n}\odot\boldsymbol{\Delta}_n^{-1}\\
\boldsymbol{\theta}_{n+1}&=\Pi_{[\boldsymbol{\ell},\mathbf{u}]}\!\left(\boldsymbol{\theta}_n-a_n\hat{\mathbf{g}}_n\right)
\end{align}
where $\Pi_{[\boldsymbol{\ell},\mathbf{u}]}$ is box projection to safety bounds.

\subsection*{PPO objective for simulation policy}

For policy $\pi_{\psi}(\Delta\boldsymbol{\theta}\mid \mathbf{o})$, use clipped objective:
\begin{equation}
\mathcal{L}^{\mathrm{CLIP}}(\psi)=
\mathbb{E}_t\left[
\min\left(
r_t(\psi)\hat{A}_t,\ 
\mathrm{clip}(r_t(\psi),1-\epsilon,1+\epsilon)\hat{A}_t
\right)
\right]
\end{equation}
with
\[
r_t(\psi)=\frac{\pi_{\psi}(a_t\mid o_t)}{\pi_{\psi_{\mathrm{old}}}(a_t\mid o_t)}
\]

The whiteboard intent is to tune architecture parameters from IMU-defined reward only. Define IMU signal vector:
\[
\mathbf{z}_t=
\begin{bmatrix}
\phi_t & \theta_t & \dot{\phi}_t & \dot{\theta}_t & \ddot{\phi}_t & \ddot{\theta}_t
\end{bmatrix}^{\top}
\]
with $\ddot{\phi}_t,\ddot{\theta}_t$ estimated by finite differences on gyro rates.

Primary IMU reward:
\[
r_t^{\mathrm{IMU}}=\exp\!\left(-\mathbf{z}_t^{\top}\mathbf{W}\mathbf{z}_t\right)-\beta_f\mathbb{I}_{\mathrm{fall},t}
\]
where $\mathbf{W}\succeq 0$ is diagonal in implementation.

Equivalent weighted linearized form (for diagnostics and ablation):
\[
r_t^{\mathrm{lin}}=
-\beta_{\phi}|\phi_t|
-\beta_{\theta}|\theta_t|
-\beta_{\dot{\phi}}|\dot{\phi}_t|
-\beta_{\dot{\theta}}|\dot{\theta}_t|
-\beta_{\ddot{\phi}}|\ddot{\phi}_t|
-\beta_{\ddot{\theta}}|\ddot{\theta}_t|
-\beta_f\mathbb{I}_{\mathrm{fall},t}.
\]

\section{Completion-Oriented Execution (No Calendar Assumption)}

\subsection*{Completion Gate 0 (Mandatory): IMU signal availability on robot}
\begin{itemize}
    \item Bring up IMU hardware and firmware path so the controller receives live measurements each fast tick
    \item Validate timestamped streams for:
    \[
    \{a_x,a_y,a_z,g_x,g_y,g_z\}_t
    \]
    with bounded drop rate and bounded jitter
    \item Estimate:
    \[
    (\phi_t,\theta_t,\dot{\phi}_t,\dot{\theta}_t,\ddot{\phi}_t,\ddot{\theta}_t)
    \]
    and verify physical plausibility under known motions
\end{itemize}
\textbf{Gate criterion 0 (must pass before any training):}
\[
\mathrm{dropout\_rate}<1\%,\quad
\mathrm{rms\_timestamp\_jitter}<0.1\Delta t,\quad
\mathrm{bias\_drift}\ \text{bounded in static test}
\]

\subsection*{Completion Block A: Architecture lock-in}
\begin{itemize}
    \item Implement $(\varphi,\varphi_l)$ oscillator formalism and unified per-leg law in \texttt{hexapod\_walking\_v3}
    \item Replace hardcoded phase functions with table-driven $(\tau_1,\tau_2,\tau_3)$ switching
    \item Define architecture interfaces explicitly:
    \[
    \mathbf{o}_t \xrightarrow{\pi_{\psi}} \Delta\boldsymbol{\theta}_t,\quad
    (\boldsymbol{\theta}_t,\varphi_t)\xrightarrow{\text{gait map}} \mathbf{q}_t,\quad
    \mathbf{q}_t\xrightarrow{\text{servo map}} \mathbf{u}^{\mu s}_t
    \]
    \item Instrument full telemetry tuple:
    \[
    \{\varphi,\varphi_l,m_l,q_l,\hat{q}_l,u^{\mu s}_{c},S_k,\texttt{mode}\}
    \]
    \item Deliver parameter-identification script for $\mu_c,d_c,u^{\min}_c,u^{\max}_c,\Delta_c^{\max}$
\end{itemize}
\textbf{Completion criterion A:}
\[
\left|\frac{T_{\mathrm{v3}}-T_{\mathrm{v2}}}{T_{\mathrm{v2}}}\right|<0.05,\quad
N_{\mathrm{clamp,v3}}\le N_{\mathrm{clamp,v2}}
\]
under matched test conditions.

\subsection*{Completion Block B: IMU reward and stabilization closure}
\begin{itemize}
    \item Integrate IMU estimator at fast tick and compute $(\phi,\theta,\dot{\phi},\dot{\theta})$
    \item Implement coupled bias law $\Delta q^{\mathrm{flx}}_{l}=s_l\Delta_{\mathrm{roll}}+r_l\Delta_{\mathrm{pitch}}$
    \item Add hysteretic mode switch $(\gamma_H,\gamma_L,M_H,M_L)$ with bounded reparameterization:
    \[
    s_{\mathrm{stride}}\!\downarrow,\ s_{\mathrm{height}}\!\uparrow,\ \omega\!\downarrow
    \]
    \item Run disturbance matrix: frontal push, lateral push, incline start, uneven contact start
\end{itemize}
\textbf{Completion criterion B:}
\[
\Pr(\mathrm{fall})_{\mathrm{reflex}} \le 0.7\,\Pr(\mathrm{fall})_{\mathrm{v2}},\quad
\bar{S}_{\mathrm{rms,reflex}}<\bar{S}_{\mathrm{rms,v2}}
\]
with identical disturbance protocol.

\textbf{Training lock:} Completion Blocks C and D are not allowed unless Gate 0 and Block B are passed.

\subsection*{Completion Block C: Parameter tuning engine}
\begin{itemize}
    \item Fix tunable architecture parameter set:
    \[
    \boldsymbol{\theta}=[K_{\phi},K_{\dot{\phi}},K_{\theta},K_{\dot{\theta}},s_{\mathrm{stride}},s_{\mathrm{height}},s_{\mathrm{freq}}]
    \]
    \item Use projected SPSA as deterministic pre-training/tuning baseline
    \item Add EEPROM persistence for best $\boldsymbol{\theta}^{\star}$ by cost $J$
    \item Produce ablation: reflex-only vs reflex+SPSA over identical terrain trials
\end{itemize}
\textbf{Completion criterion C:}
\[
\Delta J_{\mathrm{SPSA}} = J(\boldsymbol{\theta}^{\star})-J(\boldsymbol{\theta}_0) < 0
\]
with confidence interval excluding $0$ across repeated trials.

\subsection*{Completion Block D: PPO parameter tuner and transfer}
\begin{itemize}
    \item Train PPO in \texttt{sim/} with randomized contact/friction/sensor-noise/delay while keeping architecture fixed
    \item Restrict policy output to bounded parameter updates:
    \[
    \Delta\boldsymbol{\theta}_t=\pi_{\psi}(\mathbf{o}_t),\quad
    \boldsymbol{\theta}_{t+1}=\Pi_{[\ell,u]}(\boldsymbol{\theta}_{t}+\Delta\boldsymbol{\theta}_t)
    \]
    \item Validate transfer on three terrains with fixed evaluation protocol
    \item Final deliverables: architecture equations, PPO tuning equations, implementation logs, and demo script
\end{itemize}
\textbf{Completion criterion D:} PPO-assisted controller outperforms reflex+SPSA baseline on IMU reward functional:
\[
\mathcal{R}_{\mathrm{IMU}}=\frac{1}{T}\sum_{t=1}^{T} r_t^{\mathrm{IMU}}
\]
while preserving all safety constraints.

\section{Safety Envelope (Non-Negotiable)}

\begin{enumerate}
    \item \textbf{Hard pulse bounds:} never bypass \texttt{MIN\_US}/\texttt{MAX\_US}
    \item \textbf{Per-parameter bounds:} each tunable variable has fixed min/max
    \item \textbf{Slew-rate limits:} no abrupt parameter jumps between cycles
    \item \textbf{Fall detection:} if $|\phi|$ or $|\theta|$ exceeds critical angle, stop gait and enter safe mode
    \item \textbf{Recovery policy:} controlled restart sequence only
\end{enumerate}

Bounded parameter box:
\[
\ell_i \le \theta_i \le u_i,\quad
|\theta_{i,k+1}-\theta_{i,k}|\le \Delta^{\max}_{\theta_i}
\]
for every adaptive parameter $\theta_i$.

\section{Repository Integration Plan}

\begin{longtable}{>{\raggedright\arraybackslash}p{0.25\linewidth} >{\raggedright\arraybackslash}p{0.68\linewidth}}
\toprule
\textbf{Path} & \textbf{Planned role} \\
\midrule
\endhead
\texttt{firmware/hexapod\_walking\_v2/} & Frozen reference controller for regression checks \\
\texttt{firmware/hexapod\_walking\_v3/} & Runtime-parameterized gait, reflexes, adaptation hooks \\
\texttt{sim/} & PPO training, randomization, ablations, policy export \\
\texttt{host/} & Experiment runner, telemetry parser, parameter management \\
\texttt{docs/} & Benchmarks, test protocols, safety limits, milestone evidence \\
\bottomrule
\end{longtable}

\section{Practical To-Do Checklist (How To Proceed)}

\subsection*{Step 1: Bring up IMU pipeline first (hard prerequisite)}
\begin{itemize}
    \item Wire and initialize IMU in firmware; stream raw accel/gyro to logs
    \item Implement calibration routine (bias + axis sign/scale sanity checks)
    \item Add health monitors: sample frequency, dropout count, timestamp jitter
    \item Compute filtered orientation and rates:
    \[
    (\phi,\theta,\dot{\phi},\dot{\theta})
    \]
    and second derivatives via finite differences
\end{itemize}
\textbf{Artifact:} IMU health report with pass/fail on rate, jitter, and drift.

\subsection*{Step 2: Freeze baseline and create regression harness}
\begin{itemize}
    \item Duplicate current controller as immutable baseline:
    \[
    \texttt{hexapod\_walking\_v2} \rightarrow \texttt{baseline\_frozen}
    \]
    \item Create one command/script that runs fixed-duration trials and exports logs
    \item Log at minimum: phase, fall flag, clamp count, and IMU-derived state
\end{itemize}
\textbf{Artifact:} baseline log bundle + one reproducible run command.

\subsection*{Step 3: Implement architecture skeleton in v3}
\begin{itemize}
    \item Create \texttt{hexapod\_walking\_v3} with explicit pipeline:
    \[
    \mathbf{o}_t \rightarrow \Delta\boldsymbol{\theta}_t \rightarrow \mathbf{q}_t \rightarrow \mathbf{u}^{\mu s}_t
    \]
    \item Move all gait constants into a runtime parameter struct/vector
    \item Keep outputs numerically close to baseline before adding adaptation
\end{itemize}
\textbf{Artifact:} v3 firmware that reproduces v2 behavior with runtime parameters.

\subsection*{Step 4: Define IMU reward and verify observability}
\begin{itemize}
    \item Compute IMU state vector:
    \[
    \mathbf{z}_t=[\phi_t,\theta_t,\dot{\phi}_t,\dot{\theta}_t,\ddot{\phi}_t,\ddot{\theta}_t]^{\top}
    \]
    \item Implement both reward forms:
    \[
    r_t^{\mathrm{IMU}}=\exp(-\mathbf{z}_t^{\top}\mathbf{W}\mathbf{z}_t)-\beta_f\mathbb{I}_{\mathrm{fall},t},
    \quad
    r_t^{\mathrm{lin}}
    \]
    \item Verify reward traces correlate with visible stability
\end{itemize}
\textbf{Artifact:} reward debug plot and IMU-reward consistency notes.

\subsection*{Step 5: Implement leg-coupled reflex controller}
\begin{itemize}
    \item Implement per-leg flex bias:
    \[
    \Delta q^{\mathrm{flx}}_{l}=s_l\Delta_{\mathrm{roll}}+r_l\Delta_{\mathrm{pitch}}
    \]
    \item Add hysteretic cautious mode and bounded reparameterization
    \item Validate anti-phase tripod coupling remains intact
\end{itemize}
\textbf{Artifact:} disturbance-test table (baseline vs reflex).

\subsection*{Step 6: Enforce hard safety constraints everywhere}
\begin{itemize}
    \item Pulse bounds, parameter bounds, and slew bounds must all be active:
    \[
    \ell_i \le \theta_i \le u_i,\quad
    |\theta_{i,k+1}-\theta_{i,k}| \le \Delta^{\max}_{\theta_i}
    \]
    \item Add fail-safe behavior for critical tilt/fall events
    \item Record every safety intervention in logs
\end{itemize}
\textbf{Artifact:} safety validation checklist with pass/fail outcomes.

\subsection*{Step 7: Run projected SPSA as a tuning baseline}
\begin{itemize}
    \item Optimize:
    \[
    \boldsymbol{\theta}=[K_{\phi},K_{\dot{\phi}},K_{\theta},K_{\dot{\theta}},s_{\mathrm{stride}},s_{\mathrm{height}},s_{\mathrm{freq}}]
    \]
    \item Use projected update:
    \[
    \boldsymbol{\theta}_{n+1}=\Pi_{[\ell,u]}(\boldsymbol{\theta}_n-a_n\hat{\mathbf{g}}_n)
    \]
    \item Persist best $\boldsymbol{\theta}^{\star}$ to EEPROM
\end{itemize}
\textbf{Artifact:} SPSA run history + best-parameter snapshot.

\subsection*{Step 8: Train PPO as architecture-parameter tuner}
\begin{itemize}
    \item Keep gait architecture fixed; policy outputs only $\Delta\boldsymbol{\theta}$
    \item Train with randomized simulator physics/sensor delay/noise
    \item Use clipped PPO objective and IMU reward
    \item Do not train/deploy if IMU health report fails Gate 0 thresholds
\end{itemize}
\textbf{Artifact:} trained checkpoint + validation curves.

\subsection*{Step 9: Transfer PPO policy to hardware safely}
\begin{itemize}
    \item Apply policy output through projection before firmware update:
    \[
    \boldsymbol{\theta}_{t+1}=\Pi_{[\ell,u]}(\boldsymbol{\theta}_{t}+\Delta\boldsymbol{\theta}_t)
    \]
    \item Execute fixed evaluation protocol on multiple terrains
    \item Compare baseline, reflex-only, SPSA, and PPO controllers
\end{itemize}
\textbf{Artifact:} final comparison report and demo script.

\subsection*{Step 10: Lock final submission package}
\begin{itemize}
    \item Freeze code commit/tag for final architecture and policy
    \item Save reproducibility package: logs, params, scripts, plots, and PDF
    \item Ensure a single ``run-eval'' command reproduces headline metrics
\end{itemize}
\textbf{Artifact:} final reproducible release bundle.

\section{Definition of Success (Quantitative)}

Let baseline be v2 and candidate be the architecture+PPO controller. Success requires:
\begin{align}
\frac{J_{\mathrm{candidate}}-J_{\mathrm{v2}}}{J_{\mathrm{v2}}} &\le -0.25\\
\Pr(\mathrm{fall}\mid \text{disturbance set})_{\mathrm{candidate}} &\le 0.7\,
\Pr(\mathrm{fall}\mid \text{disturbance set})_{\mathrm{v2}}\\
\mathcal{R}_{\mathrm{IMU,candidate}} &> \mathcal{R}_{\mathrm{IMU,v2}}
\end{align}
plus full reproducibility of logs, parameter files, and run scripts.

\end{document}
